\documentclass{article}
\usepackage{amsmath, amssymb, geometry}
\usepackage{parskip}
\geometry{letterpaper, margin=1in}

\title{Homework 16.2 Line Integrals}
\author{Tashfeen Omran}
\date{\today}

\begin{document}

\maketitle
\tableofcontents
\newpage

\section{Part 1: Comprehensive Introduction, Context, and Prerequisites}

\subsection{Core Concepts}
A standard definite integral, $\int_a^b f(x) \, dx$, allows us to integrate a function over a straight line segment on the x-axis. A \textbf{Line Integral} (sometimes more accurately called a \textbf{path integral}) is a natural generalization of this concept. Instead of integrating along a flat, 1-dimensional axis, we integrate a function along a curve $C$ that winds through 2D or 3D space. 

There are two primary types of line integrals:
\begin{enumerate}
    \item \textbf{Line Integrals of Scalar Functions (with respect to arc length):} Denoted as $\int_C f(x, y) \, ds$ or $\int_C f(x, y, z) \, ds$. Here, we are summing up the values of a scalar function $f$ multiplied by an infinitesimally small piece of the curve's length ($ds$). Imagine a wire represented by the curve $C$. If $f(x, y)$ represents the linear mass density of the wire at point $(x, y)$, then the line integral gives the total mass of the wire.
    \item \textbf{Line Integrals of Vector Fields:} Denoted as $\int_C \mathbf{F} \cdot d\mathbf{r}$. This measures how much of a vector field $\mathbf{F}$ is aligned with the direction of the curve $C$ at every point. The most common physical interpretation is \textbf{work}: if $\mathbf{F}$ is a force field, this integral calculates the total work done by the field on a particle moving along the path $C$.
\end{enumerate}

\subsection{Intuition and Derivation}
How do we actually compute a line integral? The key is \textbf{parametrization}. We convert the curve $C$ in multi-dimensional space into a 1-dimensional interval using a parameter $t$, where $a \leq t \leq b$. 

For a scalar line integral $\int_C f(x, y) \, ds$, we parametrize $C$ as $\mathbf{r}(t) = \langle x(t), y(t) \rangle$. The differential $ds$ represents an infinitesimal arc length. By the Pythagorean theorem applied to infinitesimals:
\[ (ds)^2 = (dx)^2 + (dy)^2 \implies ds = \sqrt{\left(\frac{dx}{dt}\right)^2 + \left(\frac{dy}{dt}\right)^2} \, dt \]
This means $ds = \|\mathbf{r}'(t)\| \, dt$. Substituting this into the integral converts it back into a standard Calculus I integral with respect to $t$.

For a vector line integral $\int_C \mathbf{F} \cdot d\mathbf{r}$, the differential $d\mathbf{r}$ represents an infinitesimal displacement vector along the curve. Since $d\mathbf{r} = \mathbf{r}'(t) \, dt$, the integral becomes a standard integral of the dot product $\mathbf{F}(\mathbf{r}(t)) \cdot \mathbf{r}'(t)$ with respect to $t$.

\subsection{Historical Context and Motivation}
The concept of the line integral emerged in the early 19th century, driven heavily by the needs of classical physics. Mathematicians and physicists like George Green, Carl Friedrich Gauss, and James Clerk Maxwell needed mathematical tools to formalize the concepts of work, magnetic fields, and fluid flow. When a particle moves through a magnetic field or a gravitational field along a tortuous path, calculating the energy expended required summing up the force components continuously along that specific path. The line integral was developed specifically to solve these problems of continuous summation over arbitrary curves in space.

\subsection{Key Formulas}
\begin{itemize}
    \item \textbf{Scalar Line Integral (2D):} 
    \[ \int_C f(x, y) \, ds = \int_a^b f(x(t), y(t)) \sqrt{\left(\frac{dx}{dt}\right)^2 + \left(\frac{dy}{dt}\right)^2} \, dt \]
    \item \textbf{Scalar Line Integral (3D):} 
    \[ \int_C f(x, y, z) \, ds = \int_a^b f(x(t), y(t), z(t)) \sqrt{\left(\frac{dx}{dt}\right)^2 + \left(\frac{dy}{dt}\right)^2 + \left(\frac{dz}{dt}\right)^2} \, dt \]
    \item \textbf{Vector Line Integral (Work Integral):} 
    \[ \int_C \mathbf{F} \cdot d\mathbf{r} = \int_a^b \mathbf{F}(\mathbf{r}(t)) \cdot \mathbf{r}'(t) \, dt \]
    \item \textbf{Line Integrals with respect to coordinate variables:} 
    \[ \int_C P \, dx + Q \, dy = \int_a^b \left[ P(x(t), y(t)) x'(t) + Q(x(t), y(t)) y'(t) \right] dt \]
\end{itemize}

\subsection{Prerequisites}
To succeed in this topic, you must be proficient in:
\begin{enumerate}
    \item \textbf{Parametric Equations:} Converting Cartesian descriptions (like $y = x^2$ or a circle $x^2 + y^2 = R^2$) into parameterized forms ($x(t), y(t)$).
    \item \textbf{Vector Operations:} Computing magnitudes of vectors and taking the dot product of two vectors.
    \item \textbf{Integration Techniques:} Mastery of U-substitution and Integration by Parts is heavily required, as evaluating the final $t$-integrals often requires these methods.
\end{enumerate}

\newpage

\section{Part 2: Detailed Homework Solutions}

\subsection{Problem 1 (SCalcET9 16.2.002)}
\textbf{Question:} Evaluate the line integral, where $C$ is the given curve: $\int_C \left(\frac{x}{y}\right) \, ds$, $C: x = t^3, y = t^4, 1 \leq t \leq 4$.

\textbf{Solution:}
First, we find the derivatives of the parametric equations with respect to $t$:
\[ \frac{dx}{dt} = 3t^2, \quad \frac{dy}{dt} = 4t^3 \]
Next, we determine the arc length differential $ds$:
\[ ds = \sqrt{\left(\frac{dx}{dt}\right)^2 + \left(\frac{dy}{dt}\right)^2} \, dt = \sqrt{(3t^2)^2 + (4t^3)^2} \, dt = \sqrt{9t^4 + 16t^6} \, dt \]
Factor out $t^4$ from the expression inside the square root:
\[ ds = \sqrt{t^4(9 + 16t^2)} \, dt = t^2\sqrt{9 + 16t^2} \, dt \]
Now, substitute $x$, $y$, and $ds$ into the line integral:
\[ \int_C \left(\frac{x}{y}\right) \, ds = \int_1^4 \left(\frac{t^3}{t^4}\right) t^2\sqrt{9 + 16t^2} \, dt = \int_1^4 \left(\frac{1}{t}\right) t^2\sqrt{9 + 16t^2} \, dt = \int_1^4 t\sqrt{9 + 16t^2} \, dt \]
We evaluate this integral using $u$-substitution. Let $u = 9 + 16t^2$. Then $du = 32t \, dt$, which means $t \, dt = \frac{1}{32} du$.
Change the limits of integration:
When $t = 1$, $u = 9 + 16(1)^2 = 25$.
When $t = 4$, $u = 9 + 16(4)^2 = 9 + 256 = 265$.
\[ \int_{25}^{265} \frac{1}{32} \sqrt{u} \, du = \frac{1}{32} \left[ \frac{2}{3} u^{3/2} \right]_{25}^{265} = \frac{1}{48} \left[ 265^{3/2} - 25^{3/2} \right] \]
Since $25^{3/2} = (\sqrt{25})^3 = 5^3 = 125$, the exact final answer is:
\textbf{Final Answer:} 
\[ \frac{1}{48} (265\sqrt{265} - 125) \]

\hrule \vspace{0.5cm}

\subsection{Problem 2 (SCalcET9 16.2.003.MI)}
\textbf{Question:} Evaluate the line integral $\int_C xy^2 \, ds$, where $C$ is the right half of the circle $x^2 + y^2 = 4$ oriented counterclockwise.

\textbf{Solution:}
First, we parametrize the right half of the circle $x^2 + y^2 = 4$. The radius is $r = 2$.
\[ x(t) = 2\cos(t), \quad y(t) = 2\sin(t) \]
For the right half of the circle, the angle $t$ ranges from $-\frac{\pi}{2}$ to $\frac{\pi}{2}$.
Compute the derivatives:
\[ \frac{dx}{dt} = -2\sin(t), \quad \frac{dy}{dt} = 2\cos(t) \]
Compute $ds$:
\[ ds = \sqrt{(-2\sin(t))^2 + (2\cos(t))^2} \, dt = \sqrt{4\sin^2(t) + 4\cos^2(t)} \, dt = \sqrt{4(\sin^2(t) + \cos^2(t))} \, dt = \sqrt{4} \, dt = 2 \, dt \]
Substitute into the integral:
\[ \int_C xy^2 \, ds = \int_{-\pi/2}^{\pi/2} (2\cos(t))(2\sin(t))^2 (2 \, dt) \]
\[ = \int_{-\pi/2}^{\pi/2} (2\cos(t))(4\sin^2(t))(2) \, dt = 16 \int_{-\pi/2}^{\pi/2} \sin^2(t) \cos(t) \, dt \]
Use $u$-substitution: let $u = \sin(t)$, then $du = \cos(t) \, dt$.
When $t = -\pi/2, u = -1$. When $t = \pi/2, u = 1$.
\[ 16 \int_{-1}^1 u^2 \, du = 16 \left[ \frac{u^3}{3} \right]_{-1}^1 = 16 \left( \frac{1}{3} - \left(-\frac{1}{3}\right) \right) = 16 \left( \frac{2}{3} \right) = \frac{32}{3} \]
\textbf{Final Answer:} 
\[ \frac{32}{3} \]

\hrule \vspace{0.5cm}

\subsection{Problem 3 (SCalcET9 16.2.003.MI.SA)}
\textbf{Question:} Evaluate the line integral $\int_C xy^4 \, ds$, where $C$ is the right half of the circle $x^2 + y^2 = 9$ oriented counterclockwise.

\textbf{Solution:}
Parametrize the curve. Here, the radius is $r = \sqrt{9} = 3$.
\[ x(t) = 3\cos(t), \quad y(t) = 3\sin(t), \quad -\frac{\pi}{2} \leq t \leq \frac{\pi}{2} \]
Compute $ds$:
\[ ds = \sqrt{(-3\sin(t))^2 + (3\cos(t))^2} \, dt = \sqrt{9(\sin^2(t) + \cos^2(t))} \, dt = 3 \, dt \]
Substitute into the line integral:
\[ \int_C xy^4 \, ds = \int_{-\pi/2}^{\pi/2} (3\cos(t))(3\sin(t))^4 (3 \, dt) \]
\[ = \int_{-\pi/2}^{\pi/2} (3\cos(t))(81\sin^4(t))(3) \, dt = 729 \int_{-\pi/2}^{\pi/2} \sin^4(t) \cos(t) \, dt \]
Use $u$-substitution: let $u = \sin(t)$, then $du = \cos(t) \, dt$.
Limits change from $[-\pi/2, \pi/2]$ to $[-1, 1]$.
\[ 729 \int_{-1}^1 u^4 \, du = 729 \left[ \frac{u^5}{5} \right]_{-1}^1 = 729 \left( \frac{1}{5} - \left(-\frac{1}{5}\right) \right) = 729 \left( \frac{2}{5} \right) = \frac{1458}{5} \]
\textbf{Final Answer:} 
\[ \frac{1458}{5} \]

\hrule \vspace{0.5cm}

\subsection{Problem 4 (SCalcET9 16.2.005)}
\textbf{Question:} Evaluate the line integral $\int_C (x^6y + \sin(x)) \, dy$, where $C$ is the arc of the parabola $y = x^2$ from $(0, 0)$ to $(\pi, \pi^2)$.

\textbf{Solution:}
We can use $x$ as our parameter since the curve is given as $y$ in terms of $x$.
Let $x = t$ (which is the same as just using $x$). Then $y = x^2$.
The differential $dy$ must be expressed in terms of $dx$:
\[ dy = 2x \, dx \]
The integral ranges from $x = 0$ to $x = \pi$. Substitute $y$ and $dy$ into the integral:
\[ \int_0^\pi (x^6(x^2) + \sin(x)) (2x \, dx) = \int_0^\pi (x^8 + \sin(x))(2x) \, dx = \int_0^\pi (2x^9 + 2x\sin(x)) \, dx \]
We split this into two integrals:
\[ \int_0^\pi 2x^9 \, dx + \int_0^\pi 2x\sin(x) \, dx \]
The first integral is straightforward using the power rule:
\[ \int_0^\pi 2x^9 \, dx = \left[ \frac{2x^{10}}{10} \right]_0^\pi = \frac{\pi^{10}}{5} \]
The second integral requires \textbf{Integration by Parts}: $\int u \, dv = uv - \int v \, du$.
Let $u = 2x \implies du = 2 \, dx$.
Let $dv = \sin(x) \, dx \implies v = -\cos(x)$.
\[ \int_0^\pi 2x\sin(x) \, dx = \left[ -2x\cos(x) \right]_0^\pi - \int_0^\pi -2\cos(x) \, dx \]
\[ = \left[ -2x\cos(x) \right]_0^\pi + 2\int_0^\pi \cos(x) \, dx = \left[ -2x\cos(x) + 2\sin(x) \right]_0^\pi \]
Evaluating this from $0$ to $\pi$:
\[ (-2\pi\cos(\pi) + 2\sin(\pi)) - (-2(0)\cos(0) + 2\sin(0)) \]
\[ = (-2\pi(-1) + 0) - (0 + 0) = 2\pi \]
Combine the two results:
\textbf{Final Answer:} 
\[ \frac{\pi^{10}}{5} + 2\pi \]

\hrule \vspace{0.5cm}

\subsection{Problem 5 (SCalcET9 16.2.005.EP)}
\textbf{Question:} Consider the line integral $\int_C (x^8y + \sin(x)) \, dy$, $C$ is the arc of the parabola $y = x^2$ from $(0,0)$ to $(\pi, \pi^2)$. Use $x$ as the parameter to write parametric equations for $C$. Write an equivalent integral in only the variable $x$. Evaluate the line integral.

\textbf{Solution:}
This problem breaks down a nearly identical version of Problem 4 into steps for WebAssign.
\textbf{Part 1: Parametric Equations}
Since we use $x$ as the parameter, the equations are trivial:
\[ x = x \]
\[ y = x^2 \]
\textbf{Part 2: Equivalent Integral in $x$}
We substitute $y = x^2$ and $dy = 2x \, dx$ into the given integral:
\[ \int_C (x^8y + \sin(x)) \, dy = \int_0^\pi (x^8(x^2) + \sin(x))(2x) \, dx = \int_0^\pi (2x^{11} + 2x\sin(x)) \, dx \]
\textbf{Part 3: Evaluate}
Split into two integrals:
\[ \int_0^\pi 2x^{11} \, dx = \left[ \frac{2x^{12}}{12} \right]_0^\pi = \frac{\pi^{12}}{6} \]
The second integral is exactly the same integration by parts as Problem 4:
\[ \int_0^\pi 2x\sin(x) \, dx = \left[ -2x\cos(x) + 2\sin(x) \right]_0^\pi = 2\pi \]
\textbf{Final Answers for WebAssign Fields:} 
Parametric equations: $x = x$, $y = x^2$
Integral: $\int_0^\pi (2x^{11} + 2x\sin(x)) \, dx$
Final value: 
\[ \frac{\pi^{12}}{6} + 2\pi \]

\hrule \vspace{0.5cm}

\subsection{Problem 6 (SCalcET9 16.2.006)}
\textbf{Question:} Evaluate the line integral $\int_C e^x \, dx$, where $C$ is the arc of the curve $x = y^7$ from $(-1, -1)$ to $(1, 1)$.

\textbf{Solution:}
Here, the integral is with respect to $dx$. The curve is given as $x$ in terms of $y$. Therefore, it is most efficient to use $y$ as the parameter.
Let $y = y$, then $x = y^7$.
Find $dx$ in terms of $dy$:
\[ dx = 7y^6 \, dy \]
The limits of integration for $y$ are from $-1$ to $1$.
Substitute into the integral:
\[ \int_C e^x \, dx = \int_{-1}^1 e^{y^7} (7y^6 \, dy) \]
Use $u$-substitution: Let $u = y^7$. Then $du = 7y^6 \, dy$.
Change limits:
When $y = -1, u = (-1)^7 = -1$.
When $y = 1, u = (1)^7 = 1$.
\[ \int_{-1}^1 e^u \, du = \left[ e^u \right]_{-1}^1 = e^1 - e^{-1} \]
\textbf{Final Answer:} 
\[ e - \frac{1}{e} \]

\hrule \vspace{0.5cm}

\subsection{Problem 7 (SCalcET9 16.2.011)}
\textbf{Question:} Evaluate the line integral $\int_C xe^{yz} \, ds$, where $C$ is the line segment from $(0,0,0)$ to $(4,3,2)$.

\textbf{Solution:}
First, we must parametrize the 3D line segment. The standard parametrization for a line segment from point $P$ to point $Q$ is $\mathbf{r}(t) = (1-t)P + tQ$ for $0 \leq t \leq 1$.
Here, $P = \langle 0,0,0 \rangle$ and $Q = \langle 4,3,2 \rangle$.
\[ \mathbf{r}(t) = \langle 0,0,0 \rangle + t\langle 4-0, 3-0, 2-0 \rangle = \langle 4t, 3t, 2t \rangle, \quad 0 \leq t \leq 1 \]
So, $x = 4t$, $y = 3t$, $z = 2t$.
Compute the derivatives:
\[ x'(t) = 4, \quad y'(t) = 3, \quad z'(t) = 2 \]
Compute $ds$:
\[ ds = \sqrt{(x'(t))^2 + (y'(t))^2 + (z'(t))^2} \, dt = \sqrt{4^2 + 3^2 + 2^2} \, dt = \sqrt{16 + 9 + 4} \, dt = \sqrt{29} \, dt \]
Substitute into the line integral:
\[ \int_C xe^{yz} \, ds = \int_0^1 (4t) e^{(3t)(2t)} (\sqrt{29} \, dt) = 4\sqrt{29} \int_0^1 t e^{6t^2} \, dt \]
Use $u$-substitution: Let $u = 6t^2$. Then $du = 12t \, dt \implies t \, dt = \frac{1}{12} du$.
Limits: $t = 0 \implies u = 0$. $t = 1 \implies u = 6$.
\[ 4\sqrt{29} \int_0^6 e^u \left(\frac{1}{12} \, du\right) = \frac{4\sqrt{29}}{12} \int_0^6 e^u \, du = \frac{\sqrt{29}}{3} \left[ e^u \right]_0^6 = \frac{\sqrt{29}}{3}(e^6 - e^0) \]
\textbf{Final Answer:} 
\[ \frac{\sqrt{29}}{3}(e^6 - 1) \]

\hrule \vspace{0.5cm}

\subsection{Problem 8 (SCalcET9 16.2.011.EP)}
\textbf{Question:} Consider $\int_C xe^{yz} \, ds$, $C$ is the line segment from $(0,0,0)$ to $(4,2,3)$. Write parametric equations. Write an equivalent integral in $t$. Evaluate.

\textbf{Solution:}
Notice this is nearly identical to Problem 7, but the terminal point is $(4,2,3)$ instead of $(4,3,2)$. This swaps the $y$ and $z$ coordinates. Let's see what happens.
\textbf{Part 1: Parametric Equations}
Using $\mathbf{r}(t) = \langle 0,0,0 \rangle + t\langle 4,2,3 \rangle$:
\[ x = 4t, \quad y = 2t, \quad z = 3t, \quad 0 \leq t \leq 1 \]
\textbf{Part 2: Equivalent Integral}
Compute $ds$:
\[ ds = \sqrt{4^2 + 2^2 + 3^2} \, dt = \sqrt{16 + 4 + 9} \, dt = \sqrt{29} \, dt \]
Substitute into the integral:
\[ \int_0^1 (4t) e^{(2t)(3t)} \sqrt{29} \, dt = \int_0^1 4\sqrt{29} t e^{6t^2} \, dt \]
Because multiplication is commutative ($yz = (3t)(2t) = (2t)(3t) = 6t^2$), the resulting $t$-integral is \textit{exactly} the same as in Problem 7.
\textbf{Part 3: Evaluate}
The evaluation follows the exact same U-substitution steps as Problem 7.
\textbf{Final Answers for WebAssign Fields:}
Parametric equations: $x = 4t$, $y = 2t$, $z = 3t$
Integral: $\int_0^1 \left( 4\sqrt{29} t e^{6t^2} \right) dt$
Final value: 
\[ \frac{\sqrt{29}}{3}(e^6 - 1) \]

\hrule \vspace{0.5cm}

\subsection{Problem 9 (SCalcET9 16.2.012)}
\textbf{Question:} Evaluate the line integral $\int_C (x^2 + y^2 + z^2) \, ds$, $C: x = t, y = \cos(5t), z = \sin(5t), 0 \leq t \leq 2\pi$.

\textbf{Solution:}
First, compute the derivatives of the parametric equations:
\[ \frac{dx}{dt} = 1, \quad \frac{dy}{dt} = -5\sin(5t), \quad \frac{dz}{dt} = 5\cos(5t) \]
Compute $ds$:
\[ ds = \sqrt{(1)^2 + (-5\sin(5t))^2 + (5\cos(5t))^2} \, dt = \sqrt{1 + 25\sin^2(5t) + 25\cos^2(5t)} \, dt \]
Factor out the 25 and use the Pythagorean identity ($\sin^2(\theta) + \cos^2(\theta) = 1$):
\[ ds = \sqrt{1 + 25(\sin^2(5t) + \cos^2(5t))} \, dt = \sqrt{1 + 25(1)} \, dt = \sqrt{26} \, dt \]
Now, simplify the integrand $x^2 + y^2 + z^2$:
\[ x^2 + y^2 + z^2 = (t)^2 + (\cos(5t))^2 + (\sin(5t))^2 = t^2 + (\cos^2(5t) + \sin^2(5t)) = t^2 + 1 \]
Substitute everything into the line integral:
\[ \int_C (x^2 + y^2 + z^2) \, ds = \int_0^{2\pi} (t^2 + 1)\sqrt{26} \, dt = \sqrt{26} \int_0^{2\pi} (t^2 + 1) \, dt \]
Evaluate the definite integral:
\[ \sqrt{26} \left[ \frac{t^3}{3} + t \right]_0^{2\pi} = \sqrt{26} \left( \frac{(2\pi)^3}{3} + 2\pi - (0 + 0) \right) = \sqrt{26} \left( \frac{8\pi^3}{3} + 2\pi \right) \]
\textbf{Final Answer:} 
\[ \sqrt{26} \left( \frac{8\pi^3}{3} + 2\pi \right) \]

\hrule \vspace{0.5cm}

\subsection{Problem 10 (SCalcET9 16.2.017)}
\textbf{Question:} Evaluate the line integral $\int_C z^2 \, dx + x^2 \, dy + y^2 \, dz$, where $C$ is the line segment from $(1,0,0)$ to $(5,1,3)$.

\textbf{Solution:}
First, parametrize the line segment using $\mathbf{r}(t) = (1-t)P + tQ = P + t(Q-P)$.
$P = \langle 1,0,0 \rangle$ and $Q = \langle 5,1,3 \rangle$. The direction vector is $Q - P = \langle 4,1,3 \rangle$.
\[ \mathbf{r}(t) = \langle 1,0,0 \rangle + t\langle 4,1,3 \rangle = \langle 1+4t, t, 3t \rangle, \quad 0 \leq t \leq 1 \]
Thus:
\[ x = 1+4t \implies dx = 4 \, dt \]
\[ y = t \implies dy = 1 \, dt \]
\[ z = 3t \implies dz = 3 \, dt \]
Substitute $x, y, z$ and their differentials into the integrand:
\[ \int_C z^2 \, dx + x^2 \, dy + y^2 \, dz = \int_0^1 \left[ (3t)^2 (4 \, dt) + (1+4t)^2 (1 \, dt) + (t)^2 (3 \, dt) \right] \]
Simplify each term:
1. $z^2 \, dx = (9t^2)(4) = 36t^2$
2. $x^2 \, dy = (1 + 8t + 16t^2)(1) = 16t^2 + 8t + 1$
3. $y^2 \, dz = (t^2)(3) = 3t^2$
Combine all terms:
\[ \int_0^1 (36t^2 + 16t^2 + 8t + 1 + 3t^2) \, dt = \int_0^1 (55t^2 + 8t + 1) \, dt \]
Evaluate the definite integral:
\[ \left[ \frac{55t^3}{3} + \frac{8t^2}{2} + t \right]_0^1 = \left[ \frac{55t^3}{3} + 4t^2 + t \right]_0^1 \]
\[ = \left( \frac{55}{3} + 4 + 1 \right) - (0) = \frac{55}{3} + 5 = \frac{55}{3} + \frac{15}{3} = \frac{70}{3} \]
\textbf{Final Answer:} 
\[ \frac{70}{3} \]

\hrule \vspace{0.5cm}

\subsection{Problem 11 (SCalcET9 16.2.023.EP)}
\textbf{Question:} Consider the line integral $\int_C \mathbf{F} \cdot d\mathbf{r}$, where $\mathbf{F}(x,y,z) = \sin(x)\mathbf{i} + \cos(y)\mathbf{j} + xz\mathbf{k}$ and $\mathbf{r}(t) = t^4\mathbf{i} - t^3\mathbf{j} + t\mathbf{k}, 0 \leq t \leq 1$. Find $\mathbf{F}(\mathbf{r}(t))$, $\mathbf{r}'(t)$, their dot product, and evaluate the integral.

\textbf{Solution:}
\textbf{Part 1: Find $\mathbf{F}(\mathbf{r}(t))$}
Substitute $x = t^4$, $y = -t^3$, and $z = t$ into the vector field $\mathbf{F}$:
\[ \mathbf{F}(\mathbf{r}(t)) = \sin(t^4)\mathbf{i} + \cos(-t^3)\mathbf{j} + (t^4)(t)\mathbf{k} \]
Since cosine is an even function, $\cos(-t^3) = \cos(t^3)$. The WebAssign field accepts either, but simplified is better.
\[ \mathbf{F}(\mathbf{r}(t)) = \langle \sin(t^4), \cos(t^3), t^5 \rangle \]

\textbf{Part 2: Find $\mathbf{r}'(t)$}
Differentiate $\mathbf{r}(t) = \langle t^4, -t^3, t \rangle$ with respect to $t$:
\[ \mathbf{r}'(t) = \langle 4t^3, -3t^2, 1 \rangle \]

\textbf{Part 3: Find $\mathbf{F}(\mathbf{r}(t)) \cdot \mathbf{r}'(t)$}
Take the dot product of the vectors from Part 1 and Part 2:
\[ \langle \sin(t^4), \cos(t^3), t^5 \rangle \cdot \langle 4t^3, -3t^2, 1 \rangle = (\sin(t^4))(4t^3) + (\cos(t^3))(-3t^2) + (t^5)(1) \]
\[ = 4t^3\sin(t^4) - 3t^2\cos(t^3) + t^5 \]

\textbf{Part 4: Evaluate the Integral}
We integrate the dot product from $t = 0$ to $t = 1$:
\[ \int_0^1 (4t^3\sin(t^4) - 3t^2\cos(t^3) + t^5) \, dt \]
We can evaluate this by term-by-term integration.
For the first term, let $u = t^4 \implies du = 4t^3 \, dt$. The integral of $\sin(u)$ is $-\cos(u)$. Thus, $\int 4t^3\sin(t^4) \, dt = -\cos(t^4)$.
For the second term, let $v = t^3 \implies dv = 3t^2 \, dt$. The integral of $-\cos(v)$ is $-\sin(v)$. Thus, $\int -3t^2\cos(t^3) \, dt = -\sin(t^3)$.
For the third term, by power rule, $\int t^5 \, dt = \frac{t^6}{6}$.
Combine the antiderivatives:
\[ \left[ -\cos(t^4) - \sin(t^3) + \frac{t^6}{6} \right]_0^1 \]
Evaluate at $1$: $-\cos(1) - \sin(1) + \frac{1}{6}$
Evaluate at $0$: $-\cos(0) - \sin(0) + \frac{0}{6} = -1 - 0 + 0 = -1$
Subtract the lower limit from the upper limit:
\[ \left( -\cos(1) - \sin(1) + \frac{1}{6} \right) - (-1) = 1 + \frac{1}{6} - \cos(1) - \sin(1) = \frac{7}{6} - \cos(1) - \sin(1) \]

\textbf{Final Answers for WebAssign Fields:}
$\mathbf{F}(\mathbf{r}(t)) = \sin(t^4)\mathbf{i} + \cos(t^3)\mathbf{j} + t^5\mathbf{k}$
$\mathbf{r}'(t) = 4t^3\mathbf{i} - 3t^2\mathbf{j} + \mathbf{k}$
$\mathbf{F}(\mathbf{r}(t)) \cdot \mathbf{r}'(t) = 4t^3\sin(t^4) - 3t^2\cos(t^3) + t^5$
Value: 
\[ \frac{7}{6} - \cos(1) - \sin(1) \]

\hrule \vspace{0.5cm}

\subsection{Problem 12 (SCalcET9 16.2.023)}
\textbf{Question:} Evaluate the line integral $\int_C \mathbf{F} \cdot d\mathbf{r}$, exactly the same vector field and curve as Problem 11.

\textbf{Solution:}
This is identical to Problem 11 but presented as a single question without the required sub-steps. The calculations are precisely the same.

\textbf{Final Answer:} 
\[ \frac{7}{6} - \cos(1) - \sin(1) \]

\hrule \vspace{0.5cm}

\subsection{Problem 13 (SCalcET9 16.2.022)}
\textbf{Question:} Evaluate $\int_C \mathbf{F} \cdot d\mathbf{r}$, where $\mathbf{F}(x,y,z) = (x + y^2)\mathbf{i} + xz\mathbf{j} + (y + z)\mathbf{k}$ and $\mathbf{r}(t) = t^2\mathbf{i} + t^3\mathbf{j} - 2t\mathbf{k}, 0 \leq t \leq 2$.

\textbf{Solution:}
Extract components: $x = t^2$, $y = t^3$, $z = -2t$.
Compute $\mathbf{r}'(t)$:
\[ \mathbf{r}'(t) = \langle 2t, 3t^2, -2 \rangle \]
Substitute the parametric equations into $\mathbf{F}$:
\[ \mathbf{F}(\mathbf{r}(t)) = \langle (t^2 + (t^3)^2), (t^2)(-2t), (t^3 + (-2t)) \rangle = \langle t^2 + t^6, -2t^3, t^3 - 2t \rangle \]
Compute the dot product $\mathbf{F}(\mathbf{r}(t)) \cdot \mathbf{r}'(t)$:
\[ \langle t^6 + t^2, -2t^3, t^3 - 2t \rangle \cdot \langle 2t, 3t^2, -2 \rangle = (t^6 + t^2)(2t) + (-2t^3)(3t^2) + (t^3 - 2t)(-2) \]
\[ = (2t^7 + 2t^3) - 6t^5 - 2t^3 + 4t \]
Combine like terms (notice that $2t^3$ and $-2t^3$ cancel out):
\[ = 2t^7 - 6t^5 + 4t \]
Now, evaluate the integral from $t = 0$ to $t = 2$:
\[ \int_0^2 (2t^7 - 6t^5 + 4t) \, dt = \left[ \frac{2t^8}{8} - \frac{6t^6}{6} + \frac{4t^2}{2} \right]_0^2 = \left[ \frac{t^8}{4} - t^6 + 2t^2 \right]_0^2 \]
Evaluate at upper limit 2:
\[ \frac{2^8}{4} - 2^6 + 2(2^2) = \frac{256}{4} - 64 + 8 = 64 - 64 + 8 = 8 \]
Evaluate at lower limit 0: $0 - 0 + 0 = 0$.
\textbf{Final Answer:} 
\[ 8 \]

\hrule \vspace{0.5cm}

\subsection{Problem 14 (SCalcET9 16.2.016)}
\textbf{Question:} Evaluate the line integral $\int_C y \, dx + z \, dy + x \, dz$, where $C: x = \sqrt{t}, y = t, z = t^2, 1 \leq t \leq 9$.

\textbf{Solution:}
Compute the differentials $dx$, $dy$, and $dz$ with respect to $t$:
\[ x = t^{1/2} \implies dx = \frac{1}{2}t^{-1/2} \, dt = \frac{1}{2\sqrt{t}} \, dt \]
\[ y = t \implies dy = 1 \, dt \]
\[ z = t^2 \implies dz = 2t \, dt \]
Substitute these and the parametric equations into the integral:
\[ \int_1^9 \left[ (t)\left(\frac{1}{2\sqrt{t}} \, dt\right) + (t^2)(1 \, dt) + (\sqrt{t})(2t \, dt) \right] \]
Simplify each term:
1. $y \, dx = t \cdot \frac{1}{2}t^{-1/2} = \frac{1}{2}t^{1/2}$
2. $z \, dy = t^2$
3. $x \, dz = t^{1/2} \cdot 2t = 2t^{3/2}$
The integral becomes:
\[ \int_1^9 \left( \frac{1}{2}t^{1/2} + t^2 + 2t^{3/2} \right) dt \]
Find the antiderivative:
\[ \left[ \frac{1}{2} \frac{t^{3/2}}{3/2} + \frac{t^3}{3} + 2 \frac{t^{5/2}}{5/2} \right]_1^9 = \left[ \frac{1}{3}t^{3/2} + \frac{1}{3}t^3 + \frac{4}{5}t^{5/2} \right]_1^9 \]
Evaluate at the upper limit $t = 9$. Note that $9^{1/2} = 3$, so $9^{3/2} = 3^3 = 27$ and $9^{5/2} = 3^5 = 243$.
\[ \frac{1}{3}(27) + \frac{1}{3}(9^3) + \frac{4}{5}(243) = 9 + \frac{729}{3} + \frac{972}{5} = 9 + 243 + 194.4 = 252 + 194.4 = 446.4 = \frac{2232}{5} \]
Evaluate at the lower limit $t = 1$:
\[ \frac{1}{3}(1) + \frac{1}{3}(1) + \frac{4}{5}(1) = \frac{2}{3} + \frac{4}{5} = \frac{10}{15} + \frac{12}{15} = \frac{22}{15} \]
Subtract the lower limit evaluation from the upper limit evaluation. We use a common denominator of 15:
\[ \frac{2232}{5} - \frac{22}{15} = \frac{6696}{15} - \frac{22}{15} = \frac{6674}{15} \]
\textbf{Final Answer:} 
\[ \frac{6674}{15} \]

\newpage

\section{Part 3: In-Depth Analysis of Problems and Techniques}

\subsection{A) Problem Types and General Approach}
We can categorize the homework into four distinct problem types:

\textbf{Type 1: Direct Scalar Line Integrals with Provided Parametrizations (Problems 1, 9)}
\begin{itemize}
    \item \textit{Characteristics:} You are explicitly given $x(t), y(t), z(t)$ and the bounds for $t$. The integral contains the $ds$ term.
    \item \textit{Strategy:} Directly compute $dx/dt, dy/dt, dz/dt$. Construct $ds = \sqrt{(x')^2 + (y')^2 + (z')^2} \, dt$. Substitute $x,y,z$ into the integrand $f(x,y,z)$, multiply by the $ds$ expression, and integrate with respect to $t$.
\end{itemize}

\textbf{Type 2: Scalar Line Integrals Requiring Geometric Setup (Problems 2, 3, 7, 8)}
\begin{itemize}
    \item \textit{Characteristics:} You are given a geometric description of the curve $C$ (e.g., "right half of the circle" or "line segment from A to B").
    \item \textit{Strategy:} Step 1 is to translate the geometry into algebra by creating your own parametrization. For circles, use $x = r\cos(t), y = r\sin(t)$. For line segments from $A$ to $B$, use $\mathbf{r}(t) = A + t(B-A)$ for $0 \leq t \leq 1$. Once parameterized, it becomes a Type 1 problem.
\end{itemize}

\textbf{Type 3: Integrals with Respect to Coordinate Variables (Problems 4, 5, 6, 10, 14)}
\begin{itemize}
    \item \textit{Characteristics:} The integral does not have $ds$ or a vector dot product. Instead, it is written as $\int P \, dx + Q \, dy + R \, dz$. 
    \item \textit{Strategy:} These are often the easiest to set up. If the curve is given as $y = f(x)$, let $x = t$ (so $dx = dt$), and $dy = f'(x) \, dx$. Substitute every variable and differential to convert the entire expression into a single variable (like $x$, $y$, or $t$) and integrate normally.
\end{itemize}

\textbf{Type 4: Vector Field Line Integrals (Problems 11, 12, 13)}
\begin{itemize}
    \item \textit{Characteristics:} Written as $\int_C \mathbf{F} \cdot d\mathbf{r}$. You are given a vector field $\mathbf{F}$ and a parametric curve $\mathbf{r}(t)$.
    \item \textit{Strategy:} Follow the exact sequence: (1) Find $\mathbf{F}(\mathbf{r}(t))$ by substituting the $x,y,z$ components of $\mathbf{r}(t)$ into $\mathbf{F}$. (2) Find the velocity vector $\mathbf{r}'(t)$. (3) Compute the algebraic dot product of these two vectors to yield a scalar function of $t$. (4) Integrate this scalar function.
\end{itemize}

\subsection{B) Key Algebraic and Calculus Manipulations}
\begin{itemize}
    \item \textbf{Pythagorean Algebraic Simplification:} In Problem 9, calculating $ds$ resulted in $\sqrt{1 + 25\sin^2(5t) + 25\cos^2(5t)} \, dt$. Recognizing that you can factor out the $25$ to create the identity $\sin^2(\theta) + \cos^2(\theta) = 1$ is an essential trick. It magically removes the trigonometric functions from under the radical, making the integral solvable.
    \item \textbf{U-Substitution for 'Inner' Functions:} Problems 1, 6, and 7 heavily relied on recognizing a function and its derivative. In Problem 7 ($\int 4\sqrt{29} t e^{6t^2} \, dt$), recognizing that the derivative of the exponent $6t^2$ is $12t$ allowed for a clean u-substitution to eliminate the $t$ outside the exponential.
    \item \textbf{Integration by Parts for Mixed Types:} Problems 4 and 5 required integrating $2x\sin(x)$. Whenever you have a polynomial multiplied by a trigonometric or exponential function, it signals the need for Integration by Parts ($\int u \, dv = uv - \int v \, du$). The "trick" is to let $u$ be the polynomial so that taking its derivative reduces its degree.
    \item \textbf{Symmetry/Commutativity Trickery:} Problems 7 and 8 demonstrated a WebAssign trick. By swapping the $y$ and $z$ terminal coordinates from $(4,3,2)$ to $(4,2,3)$, the required integral changed from integrating $e^{(3t)(2t)}$ to $e^{(2t)(3t)}$. Because multiplication is commutative, the resulting math is 100\% identical. Recognizing this saves minutes of redundant calculation.
\end{itemize}

\newpage

\section{Part 4: "Cheatsheet" and Tips for Success}

\subsection{Must-Know Formulas Summary}
\begin{itemize}
    \item \textbf{Arc Length Differential:} $ds = \|\mathbf{r}'(t)\| \, dt = \sqrt{(x')^2 + (y')^2 + (z')^2} \, dt$
    \item \textbf{Line Segment Parametrization:} To go from point $P$ to $Q$: $\mathbf{r}(t) = (1-t)P + tQ$, where $0 \leq t \leq 1$.
    \item \textbf{Circle Parametrization:} $x(t) = R\cos(t)$, $y(t) = R\sin(t)$. Standard counterclockwise orientation.
    \item \textbf{Work Integral Definition:} $\int_C \mathbf{F} \cdot d\mathbf{r} = \int_a^b \left[ \mathbf{F}(\mathbf{r}(t)) \cdot \mathbf{r}'(t) \right] dt$
\end{itemize}

\subsection{Speed Hacks and Tricks}
\begin{itemize}
    \item \textbf{The "Function as Parameter" Hack:} If a curve is given as $y = f(x)$, don't introduce a new variable $t$. Just use $x$ as your parameter. $x = x$, $dx = dx$, $y = f(x)$, $dy = f'(x) \, dx$. This instantly sets up integrals of the form $\int P \, dx + Q \, dy$. 
    \item \textbf{Dot Product Check:} Before integrating $\mathbf{F}(\mathbf{r}(t)) \cdot \mathbf{r}'(t)$, quickly scan the terms. If you see terms that are exact opposites (like $2t^3$ and $-2t^3$ in Problem 13), cancel them immediately before integrating to save time and reduce arithmetic errors.
    \item \textbf{Even/Odd Function Simplification:} If you have $\cos(-x)$ resulting from substituting a negative parameter, immediately write it as $\cos(x)$ because cosine is an even function. If you have $\sin(-x)$, pull the negative out to make it $-\sin(x)$.
\end{itemize}

\subsection{Common Traps and Pitfalls}
\begin{itemize}
    \item \textbf{Forgetting to substitute $x, y, z$ into $\mathbf{F}$:} The most common student error in vector line integrals is leaving $\mathbf{F}$ in terms of $x, y, z$ while trying to integrate with respect to $t$. You \textit{must} plug your $\mathbf{r}(t)$ components into $\mathbf{F}$ first!
    \item \textbf{Messing up $ds$ vs $d\mathbf{r}$:} A scalar integral uses $ds$ (which involves a square root magnitude). A vector integral uses $d\mathbf{r}$ (which involves a dot product). Do not put a square root magnitude inside a vector work integral.
    \item \textbf{Reversing the Orientation:} Integrating from $A$ to $B$ is not the same as integrating from $B$ to $A$ for vector integrals. If you write your parametrization backward, your final answer will have the wrong sign. $\int_{-C} \mathbf{F} \cdot d\mathbf{r} = -\int_C \mathbf{F} \cdot d\mathbf{r}$. (Note: This is \textit{not} true for scalar line integrals, which are path-direction independent!).
\end{itemize}

\newpage

\section{Part 5: Conceptual Synthesis and The "Big Picture"}

\subsection{A) Thematic Connections}
The core theme of a line integral is \textbf{"accumulation along a path."} This represents a massive leap in how we view integration.
In Calculus I, integration was rigidly defined over straight intervals on the x-axis (Riemann sums of rectangles). The "base" of every rectangle was just $dx$.
With line integrals, we bend that x-axis into a flexible wire that can curve through space. The fundamental theme remains exactly the same—we are summing up an infinite number of infinitesimally small pieces. However, the "base" of our geometric pieces is now an arc length $ds$ or a displacement vector $d\mathbf{r}$.
This theme appears identically in \textbf{Surface Integrals} (which you will study soon), where instead of bending a 1D line through space, we bend a 2D plane into a curved surface and integrate over surface area patches $dS$. 

\subsection{B) Forward and Backward Links}
\textbf{Backward Links:} Line integrals are the ultimate synthesis of Calculus I, II, and early Calculus III. They require the parametric equations learned in Calc II, the vector operations learned at the start of Calc III, and the fundamental integration techniques (U-sub, Parts) from Calc I/II. 

\textbf{Forward Links:} Line integrals are the bedrock for the most important theorems in multivariable calculus:
\begin{itemize}
    \item \textbf{The Fundamental Theorem of Line Integrals:} Shows how line integrals of conservative fields relate to potential functions (analogous to the Fundamental Theorem of Calculus).
    \item \textbf{Green's Theorem:} Relates a line integral around a closed boundary curve to a double integral over the region it encloses.
    \item \textbf{Stokes' Theorem:} The 3D extension of Green's theorem, relating the line integral around the boundary of a 3D surface to the surface integral of the "curl" of the vector field. 
\end{itemize}

\newpage

\section{Part 6: Real-World Application and Modeling}

\subsection{A) Concrete Scenarios in Quantitative Finance and Statistics}
While traditionally used for calculating physical "work," line integrals are deeply integrated into advanced finance, specifically in stochastic continuous-time finance.

\textbf{Scenario 1: Path-Dependent Derivative Pricing (Asian Options)}
An Asian option is a financial derivative whose payoff depends on the average price of the underlying asset over a certain period, not just its price at maturity. If we view the asset's price history as a curve $C$ parametrized by time $t$ in a price-time plane, the accumulation of the continuous average is mathematically formulated using a line integral along the asset's specific price path.

\textbf{Scenario 2: Computing the "Work" of Rebalancing a Portfolio}
Imagine a portfolio containing two assets, $X$ and $Y$. The state of the portfolio over time forms a parametric curve $C$ in the $XY$-plane. If there are continuous transaction costs that depend on the current prices of the assets and the rate at which you buy/sell them, the total transaction cost over a quarter is the line integral of a "cost vector field" along the portfolio's rebalancing path.

\textbf{Scenario 3: The It\^o Integral and Stochastic Calculus}
In the Black-Scholes-Merton model, stock prices follow a random walk (Brownian motion). To calculate the value of a portfolio over time, quants use the It\^o Integral, denoted $\int_0^T \Delta \, dW_t$. This is conceptually a line integral, but instead of integrating over a smooth, deterministic differential $d\mathbf{r}$, we integrate over a stochastic differential $d\mathbf{W}_t$ (the increment of Brownian motion). Line integrals form the cognitive stepping stone to this extremely advanced financial mathematics.

\subsection{B) Model Problem Setup: Portfolio Rebalancing Costs}
Let's formalize \textbf{Scenario 2}. 

\textbf{Variables and Definitions:}
Let $x(t)$ be the number of shares held of Stock A at time $t$.
Let $y(t)$ be the number of shares held of Stock B at time $t$.
The portfolio evolution over time $t$ from $t=0$ to $t=1$ (one year) defines a curve $C$ parametrized by $\mathbf{r}(t) = \langle x(t), y(t) \rangle$.

Suppose the broker charges a transaction fee that is proportional to the size of the trade, but the fee rate fluctuates based on market liquidity. We can model the marginal cost of trading Stock A as $P(x,y) = 0.01x + 0.05y$ and Stock B as $Q(x,y) = 0.05x + 0.02y$. This creates a "Cost Vector Field": 
\[ \mathbf{F}(x,y) = \langle P(x,y), Q(x,y) \rangle \]

\textbf{The Mathematical Model:}
To find the Total Transaction Cost ($TC$) accumulated by continuously rebalancing the portfolio along path $C$ over the year, a quantitative analyst must compute the vector line integral of the cost field along the trading path:
\[ TC = \int_C \mathbf{F} \cdot d\mathbf{r} = \int_C (0.01x + 0.05y) \, dx + (0.05x + 0.02y) \, dy \]
By plugging in the specific trading strategies $x(t)$ and $y(t)$, the firm can compute exactly how much money will be lost to friction costs over the year.

\newpage

\section{Part 7: Common Variations and Untested Concepts}

The homework set provided primarily focused on the mechanics of computing line integrals algebraically. However, a major concept in standard Calculus III curricula was entirely absent: \textbf{Conservative Vector Fields and the Fundamental Theorem of Line Integrals (FTLI)}.

\textbf{The Missing Concept: Path Independence}
Normally, the value of a line integral $\int_C \mathbf{F} \cdot d\mathbf{r}$ depends on the specific path taken from point A to point B. However, if the vector field $\mathbf{F}$ is \textit{conservative} (meaning it is the gradient of some scalar potential function $f$, so $\nabla f = \mathbf{F}$), then the line integral is \textbf{path independent}. 

According to the Fundamental Theorem of Line Integrals:
\[ \int_C \mathbf{F} \cdot d\mathbf{r} = \int_C \nabla f \cdot d\mathbf{r} = f(\mathbf{r}(b)) - f(\mathbf{r}(a)) \]
This means you don't need to parametrize the curve at all! You just evaluate the potential function $f$ at the end point and subtract its value at the starting point.

\textbf{Worked Example of the Missing Concept:}
\textit{Problem:} Evaluate $\int_C \mathbf{F} \cdot d\mathbf{r}$, where $\mathbf{F}(x,y) = \langle 2xy, x^2 \rangle$, and $C$ is a wildly complicated spiraling path given by $\mathbf{r}(t) = \langle t\cos^5(t), e^{\sin(t)} \rangle$ from $t = 0$ to $t = \pi$.

\textit{Standard approach (impossible here):} Plugging $\mathbf{r}(t)$ into $\mathbf{F}$ and multiplying by $\mathbf{r}'(t)$ would create an algebraically unsolvable nightmare.

\textit{FTLI Approach:} 
1. Check if $\mathbf{F}$ is conservative. Is $\frac{\partial}{\partial y}(2xy) = \frac{\partial}{\partial x}(x^2)$? Yes, $2x = 2x$. 
2. Find the potential function $f(x,y)$ such that $\nabla f = \mathbf{F}$. 
   Integrate the $x$-component with respect to $x$: $f(x,y) = \int 2xy \, dx = x^2y + g(y)$.
   Differentiate with respect to $y$ and set equal to the $y$-component: $x^2 + g'(y) = x^2 \implies g'(y) = 0 \implies g(y) = K$.
   So, $f(x,y) = x^2y$.
3. Find the start and end points of the curve.
   Start point ($t=0$): $\mathbf{r}(0) = \langle 0\cos^5(0), e^{\sin(0)} \rangle = \langle 0, e^0 \rangle = (0, 1)$.
   End point ($t=\pi$): $\mathbf{r}(\pi) = \langle \pi\cos^5(\pi), e^{\sin(\pi)} \rangle = \langle -\pi, 1 \rangle$.
4. Apply the Theorem:
   \[ \int_C \mathbf{F} \cdot d\mathbf{r} = f(-\pi, 1) - f(0, 1) = ((-\pi)^2(1)) - ((0)^2(1)) = \pi^2 \]
By understanding this untested concept, you can bypass impossible integrals instantly.

\newpage

\section{Part 8: Advanced Diagnostic Testing: "Find the Flaw"}

Below are 5 solved problems. The solutions look correct but contain a subtle, critical flaw. Your task is to act as the grader: find the error, explain why it is wrong in one sentence, and provide the correct step.

\subsection{Flawed Problem 1}
\textbf{Problem:} Evaluate the scalar line integral $\int_C 2xy \, ds$, where $C$ is the line segment from $(0,0)$ to $(1,2)$ defined by $\mathbf{r}(t) = \langle t, 2t \rangle$ for $0 \leq t \leq 1$.
\textbf{Flawed Solution:}
1. Extract $x(t) = t$ and $y(t) = 2t$.
2. Substitute into the integral: $\int_0^1 2(t)(2t) \, dt$.
3. Simplify: $\int_0^1 4t^2 \, dt$.
4. Integrate: $\left[ \frac{4t^3}{3} \right]_0^1 = \frac{4}{3}$.
\textbf{Correct Answer:} \textit{[Your job to find the flaw!]}

\subsection{Flawed Problem 2}
\textbf{Problem:} Evaluate $\int_C y \, dx + x \, dy$, where $C$ is the curve $y = x^3$ from $(1,1)$ to $(2,8)$.
\textbf{Flawed Solution:}
1. Let $x = t$, so $dx = dt$. Then $y = t^3$, so $dy = 3t^2 \, dt$.
2. Substitute into the integral: $\int_1^2 (t^3)(dt) + (t)(3t^2) \, dt$.
3. Simplify: $\int_1^2 (t^3 + 3t^3) \, dt = \int_1^2 4t^3 \, dt$.
4. Integrate: $\left[ t^4 \right]_1^2 = 2^4 - 1^4 = 16 - 1 = 15$.
5. Since $y=x^3$ was a curve backwards from the origin, we flip the sign of the final answer: $-15$.
\textbf{Correct Answer:} \textit{[Your job to find the flaw!]}

\subsection{Flawed Problem 3}
\textbf{Problem:} Evaluate the vector line integral $\int_C \mathbf{F} \cdot d\mathbf{r}$ for $\mathbf{F} = \langle y, x \rangle$ and $\mathbf{r}(t) = \langle t, t^2 \rangle$ for $0 \leq t \leq 2$.
\textbf{Flawed Solution:}
1. Find $\mathbf{r}'(t) = \langle 1, 2t \rangle$.
2. The dot product is $\mathbf{F} \cdot \mathbf{r}'(t) = (y)(1) + (x)(2t) = y + 2xt$.
3. Integrate: $\int_0^2 (y + 2xt) \, dt = \left[ yt + xt^2 \right]_0^2 = 2y + 4x$.
\textbf{Correct Answer:} \textit{[Your job to find the flaw!]}

\subsection{Flawed Problem 4}
\textbf{Problem:} Evaluate the scalar line integral $\int_C x^2 \, ds$ where $C$ is the left half of the circle $x^2 + y^2 = 4$ oriented from $(0,2)$ down to $(0,-2)$.
\textbf{Flawed Solution:}
1. Parametrize the circle: $x(t) = 2\cos(t), y(t) = 2\sin(t)$.
2. To get the left half from top to bottom, let $t$ range from $\frac{\pi}{2}$ to $\frac{3\pi}{2}$.
3. Because the curve is going "backwards" (downward), scalar line integrals flip sign when path is reversed. 
4. The integral is $-\int_{\pi/2}^{3\pi/2} (2\cos(t))^2 \sqrt{(-2\sin t)^2 + (2\cos t)^2} \, dt = -8 \int_{\pi/2}^{3\pi/2} \cos^2(t) \, dt = -4\pi$.
\textbf{Correct Answer:} \textit{[Your job to find the flaw!]}

\subsection{Flawed Problem 5}
\textbf{Problem:} Evaluate $\int_C x \, dx + y \, dy$ where $C$ is the curve $y = \sin(x)$ from $(0,0)$ to $(\pi, 0)$.
\textbf{Flawed Solution:}
1. Let $x = t$, $y = \sin(t)$.
2. $dx = dt$, $dy = \cos(t) \, dt$.
3. Substitute: $\int_0^\pi t \, dt + \sin(t)\cos(t) \, dt$.
4. Integrate $t$: $\frac{t^2}{2}$. Integrate $\sin(t)\cos(t)$ using u-sub $u=\cos(t), du=-\sin(t) dt$: $-\frac{\cos^2(t)}{2}$.
5. Evaluate $\left[ \frac{t^2}{2} - \frac{\cos^2(t)}{2} \right]_0^\pi = (\frac{\pi^2}{2} - \frac{1}{2}) - (0 - \frac{1}{2}) = \frac{\pi^2}{2} - 1$.
\textbf{Correct Answer:} \textit{[Your job to find the flaw!]}

\end{document}